\documentclass[12pt]{article}

\usepackage[a4paper,margin=2.5cm]{geometry}

\usepackage[onehalfspacing]{setspace}

\usepackage{amsmath}
\usepackage{amssymb}
\usepackage{xcolor}
\usepackage{mathtools}
\usepackage{amsthm}
\usepackage{tikz-cd}
\newtheorem{definition}{Definition}
\newtheorem{example}{Example}

\newenvironment{solution}
{\begin{proof}[Solution]}
{\end{proof}}

\newenvironment{Remark}
{\begin{proof}[Remark]}
{\end{proof}}

\newenvironment{Goal}
{\begin{proof}[Goal]}
{\end{proof}}

\setlength{\parindent}{0pt}

\begin{document}

\begin{titlepage}
    \centering

    \vspace*{\fill}

    {\Huge Linear Transformations IV}

    \vspace{1cm}

    {\Large Jin-Ren Ryan Wang\par}

    \vspace{0.5cm}

    {\large August 9, 2026\par}

    \vspace*{\fill}
\end{titlepage}

\tableofcontents

\newpage

\section{Recap on elementary matrices}

\subsection{Definition}

\textbf{Definition 7.1.1.}
We identify elementary row matrices and elementary column matrices as follows:
\[
E_{i,j}^{\mathrm{I}}
    = E_{\mathrm{I}}^{\,i,j}, \qquad
E_{\lambda,i}^{\mathrm{II}}
    = E_{\mathrm{II}}^{\,\lambda,i}, \qquad
E_{\mu,i,j}^{\mathrm{III}}
    = E_{\mathrm{III}}^{\,\mu,i,j}.
\]

From now on, we simply refer to them as \emph{elementary matrices}.

\subsection{Theorem}

\textbf{Theorem 7.1.1}
Every invertible matrix can be written as a product of elementary matrices.

\begin{proof}
Let $A$ be an invertible matrix. Since $A$ is invertible,
\begin{enumerate}
    \item $A$ is an $n\times n$ square matrix;
    \item $\operatorname{rank}(A)=n$.
\end{enumerate}

By the QAP Theorem, there exist elementary matrices $\textcolor{red}{E_1,E_2,\ldots,E_k},\quad G_1,G_2,\ldots,G_\ell$,
s.t. $\textcolor{red}{(E_1E_2\cdots E_k)}\,\textcolor{blue}{A}\,(G_1G_2\cdots G_\ell)= \underbrace{I_n}_{\textcolor{red}{\text{SNF of $A$}}}$ .

Hence, $A=\textcolor{red}{(E_1E_2\cdots E_k)^{-1}}\textcolor{blue}{I_n}(G_1G_2\cdots G_\ell)^{-1}$,

$\implies$ $A=\textcolor{red}{E_k^{-1}\cdots E_2^{-1}E_1^{-1}G_\ell^{-1}\cdots G_2^{-1}G_1^{-1}}$.

Since the inverse of every elementary matrix is again an elementary matrix, $A$ is a product of elementary matrices.
\end{proof}

\section{Augmented matrices}

\textbf{Definition}

\textbf{Definition 7.2.1.}
Let $A\in M_{m\times p}(F), \qquad B\in M_{m\times q}(F)$. The \emph{augmented matrix} $(A\mid B)$ is an
$m\times (p+q)$ matrix for which the first $p$ columns give the matrix $A$ and the last $q$ columns give the matrix $B$.

\textbf{Note. }
The vertical line is placed to separate the matrices $A$ and $B$ for convenience only.

Let
\[
A=
\begin{pmatrix}
1 & 2\\
3 & 4
\end{pmatrix}.
\]

Then
\[
(A\mid I_2)=
\left(
\begin{array}{cc|cc}
1 & 2 & 1 & 0\\
3 & 4 & 0 & 1
\end{array}
\right).
\]

\newpage

\section{Application: Finding inverse of a matrix}

Let $A$ be an $(n\times n)$ invertible matrix. Consider the augmented matrix
$(A\mid I_n)$ formed by $A$ and $I_n$. The key observation is that
\[
\textcolor{red}{A^{-1}}(A\mid I_n)=(I_n\mid \textcolor{red}{A^{-1}}).
\]

By Theorem 7.1.1, $A^{-1}$, as an invertible matrix, can be written as a
product of elementary matrices. Therefore,
\[
\underbrace{\textcolor{red}{(E_1E_2\cdots E_k)}}_{\textcolor{red}{A^{-1}}}(A\mid I_n)=(I_n\mid A^{-1}),
\]
where
\[
A^{-1}=E_1E_2\cdots E_k.
\]

Since multiplying elementary matrices from the left corresponds to
conducting \emph{row operations}, the above operations give us a concrete
strategy to compute the inverse of a given matrix $A$, which we record below.

\subsection{Theorem}

\textbf{Theorem 7.3.1.}
Let $A$ be an $(n\times n)$ invertible matrix. If we conduct row operations
to transform the augmented matrix
\[
(A\mid I_n)\xrightarrow{\text{row operations}}(I_n\mid B),
\]
Then $B=A^{-1}$.
\subsection{Example}
\subsubsection{Example 1}
Compute the inverse of the matrix
\[
A=
\begin{pmatrix}
1 & 2 & 1\\
0 & 1 & 3\\
2 & 4 & -1
\end{pmatrix}.
\]

\begin{solution}
We aim to transform the augmented matrix $(A\mid I_3)$ into the form
$(I_3\mid B)$ by row operations.

\begin{align*}
(A\mid I_3)
&=
\left(
\begin{array}{ccc|ccc}
1 & 2 & 1  & 1 & 0 & 0\\
0 & 1 & 3  & 0 & 1 & 0\\
2 & 4 & -1 & 0 & 0 & 1
\end{array}
\right) \\[0.5em]
&\longrightarrow
\left(
\begin{array}{ccc|ccc}
1 & 2 & 1  & 1  & 0 & 0\\
0 & 1 & 3  & 0  & 1 & 0\\
0 & 0 & -3 & -2 & 0 & 1
\end{array}
\right) \\[0.5em]
&\longrightarrow
\left(
\begin{array}{ccc|ccc}
1 & 2 & 1 & 1           & 0 & 0\\
0 & 1 & 3 & 0           & 1 & 0\\
0 & 0 & 1 & \frac{2}{3} & 0 & -\frac{1}{3}
\end{array}
\right) \\[0.5em]
&\longrightarrow
\left(
\begin{array}{ccc|ccc}
1 & 2 & 0 & \frac{1}{3}  & 0 & \frac{1}{3}\\
0 & 1 & 0 & -2            & 1 & 1\\
0 & 0 & 1 & \frac{2}{3}   & 0 & -\frac{1}{3}
\end{array}
\right) \\[0.5em]
&\longrightarrow
\left(
\begin{array}{ccc|ccc}
1 & 0 & 0 & \frac{13}{3} & -2 & -\frac{5}{3}\\
0 & 1 & 0 & -2           & 1  & 1\\
0 & 0 & 1 & \frac{2}{3}  & 0  & -\frac{1}{3}
\end{array}
\right).
\end{align*}

Thus,
\[
\boxed{
A^{-1}
=
\begin{pmatrix}
\frac{13}{3} & -2 & -\frac{5}{3}\\
-2 & 1 & 1\\
\frac{2}{3} & 0 & -\frac{1}{3}
\end{pmatrix}
}.
\]
\end{solution}

\subsubsection{Example 2}
Consider the function $T:\mathbb{R}[x]_{\leq 2}\to\mathbb{R}[x]_{\leq 2}$
defined by
\[
T(p(x))=p(x)+p'(x)+p''(x).
\]
Let $B=\{1,x,x^2\}$ be a basis of $\mathbb{R}[x]_{\leq 2}$.

\begin{enumerate}
    \item[(a)] Prove that $T$ is a linear transformation.
    \item[(b)] Find $[T]^B_B$. Is $T$ an isomorphism?
    \item[(c)] Write down the inverse function $T^{-1}$.
\end{enumerate}

\begin{solution}
\begin{enumerate}
\item[(a)]
For any $p(x),q(x)\in\mathbb{R}[x]_{\leq 2}$,
\begin{align*}
T(p(x)+q(x))
&=(p(x)+q(x))+(p(x)+q(x))'+(p(x)+q(x))''\\
&=p(x)+q(x)+p'(x)+q'(x)+p''(x)+q''(x)\\
&=\bigl(p(x)+p'(x)+p''(x)\bigr)
 +\bigl(q(x)+q'(x)+q''(x)\bigr)\\
&=T(p(x))+T(q(x)).
\end{align*}

Also, for any $\alpha\in\mathbb{R}$,
\begin{align*}
T(\alpha p(x))
&=\alpha p(x)+(\alpha p(x))'+(\alpha p(x))''\\
&=\alpha p(x)+\alpha p'(x)+\alpha p''(x)\\
&=\alpha\bigl(p(x)+p'(x)+p''(x)\bigr)\\
&=\alpha T(p(x)).
\end{align*}

Therefore, $T$ is a linear transformation.

\item[(b)]
We have
\begin{align*}
T(\textcolor{red}{1})
&=1
 =1\cdot1+0\cdot x+0\cdot x^2,\\
T(\textcolor{red}{x})
&=x+1
 =1\cdot1+1\cdot x+0\cdot x^2,\\
T(\textcolor{red}{x^2})
&=x^2+2x+2
 =2\cdot1+2\cdot x+1\cdot x^2.
\end{align*}

So,
$
[T]^B_B=
\begin{pmatrix}
1 & 1 & 2\\
0 & 1 & 2\\
0 & 0 & 1
\end{pmatrix}.
$ is already in REF, and
\[
\operatorname{rank}[T]^B_B=3=\dim(\mathbb{R}[x]_{\leq2}) \implies \text{surjective}
\]
\[
\operatorname{nullity}[T]^B_B=\underbrace{\textcolor{red}{3}}_{\textcolor{red}{\text{dim of domain}}}-\underbrace{3}_{\text{rank}}=0 \implies \text{injective}
\]
so $T$ is an isomorphism.

\item[(c)]
We first compute the inverse matrix. Note: $[T^{-1}]_B^B = ([T]_B^B)^{-1}$

Consider
\[
\left(
\begin{array}{ccc|ccc}
1&1&2&1&0&0\\
0&1&2&0&1&0\\
0&0&1&0&0&1
\end{array}
\right)
\longrightarrow
\left(
\begin{array}{ccc|ccc}
1&1&0&1&0&-2\\
0&1&0&0&1&-2\\
0&0&1&0&0&1
\end{array}
\right)
\]
and then
\[
\longrightarrow
\left(
\begin{array}{ccc|ccc}
1&0&0&1&-1&0\\
0&1&0&0&1&-2\\
0&0&1&0&0&1
\end{array}
\right).
\]

Therefore,
\[
[T^{-1}]^B_B
=
\left([T]^B_B\right)^{-1}
=
\begin{pmatrix}
1&-1&0\\
0&1&-2\\
0&0&1
\end{pmatrix}.
\]

For
\[
p(x)=a+bx+cx^2,
\]
we have
\[
\begin{pmatrix}
1&-1&0\\
0&1&-2\\
0&0&1
\end{pmatrix}
\begin{pmatrix}
a\\b\\c
\end{pmatrix}
=
\begin{pmatrix}
a-b\\
b-2c\\
c
\end{pmatrix}.
\]

Translating back to $\mathbb{R}[x]_{\leq2}$,
\[
\boxed{
T^{-1}(a+bx+cx^2)
=
(a-b)+(b-2c)x+cx^2
}.
\]

\end{enumerate}
\end{solution}

\section{Application: System of linear equations}

Historically, matrices arise when one attempts to solve systems of linear
equations. For example, consider the system
\[
\begin{cases}
x+2y=3,\\
3x+4y=1.
\end{cases}
\]

A condensed way to write the above system is
$\left(
\begin{array}{cc|c}
1 & 2 & 3\\
3 & 4 & 1
\end{array}
\right).$ This is called the \emph{augmented matrix} for the above system of equations.

\subsection{Definition}

\textbf{Definition 7.4.1.}
Let $(A\mid \mathbf{b})$ be the augmented matrix associated with a system
of linear equations $A\mathbf{x}=\mathbf{b}$. We call the system of equations
\begin{itemize}
    \item \emph{homogeneous} if $\mathbf{b}=\mathbf{0}$;
    \item \emph{inhomogeneous} if $\mathbf{b}\neq\mathbf{0}$.
\end{itemize}

\subsection{Theorem}

\textbf{Theorem 7.4.1.}
The solution set of a homogeneous system $(A\mid \mathbf{0})$ coincides
with the subspace $\ker(A)$. In particular,
\[
\underbrace{dim(\ker(A))}_{\operatorname{nullity}(A)}=(\text{number of columns of }A)-\operatorname{rank}(A).
\]

\begin{proof}
A vector $\mathbf{x}$ is a solution of $(A\mid\mathbf{0})$ $\Leftrightarrow$ $A\mathbf{x}=\mathbf{0}$

$\implies$ $(A\mid\mathbf{0})$ $\Leftrightarrow$ $\mathbf{x}\in\ker(A)$.

Moreover, by the rank-nullity theorem,
\[
\dim(\ker(A))=\operatorname{nullity}(A)=(\text{number of columns of }A)-\operatorname{rank}(A).
\]
\end{proof}

\textbf{Theorem 7.4.2.}
Let $A\in M_{m\times n}(F)$. If $m<n$, then the homogeneous system
$(A\mid\mathbf{0})$ has a non-zero solution.

\begin{proof}
Since $A\in M_{m\times n}(F)$ represents a linear map $F^n \longrightarrow F^m$,
we have $\operatorname{im}(A)\subseteq F^m$. Therefore, $\operatorname{rank}(A)\leq m$.
By the rank-nullity theorem,
\begin{align*}
\dim(\ker(A))
&=(\text{number of columns of }A)-\operatorname{rank}(A)\\
&=n-\operatorname{rank}(A)\\
&\geq n-m \textcolor{red}{\text{ $\geq 1$}} \implies \textcolor{red}{\text{$n \geq m$}}
\end{align*}
Hence, $\ker(A)\neq\{\mathbf{0}\}$. Therefore, there $\exists$ some non-zero vector $\mathbf{x}\in\ker(A)$
s.t. $A\mathbf{x}=\mathbf{0}$.

Thus, the homogeneous system $(A\mid\mathbf{0})$ has a non-zero solution.
\end{proof}

\textbf{Theorem 7.4.3.}
Let $(A\mid\mathbf{b})$ be the augmented matrix of a general system of
equations. Then the system is \emph{consistent} (has a solution) iff.
\[
\operatorname{rank}(A)=\operatorname{rank}(A\mid\mathbf{b}).
\]

\begin{proof}
Let $A=
\begin{pmatrix}
\mathbf{c}_1 & \mathbf{c}_2 & \cdots & \mathbf{c}_n
\end{pmatrix}$.

$(\Rightarrow)$

Suppose $\mathbf{x}=
\begin{pmatrix}
x_1\\
x_2\\
\vdots\\
x_n
\end{pmatrix}$
is a solution of $A\mathbf{x}=\mathbf{b}$.

$\implies$ $x_1\mathbf{c}_1+x_2\mathbf{c}_2+\cdots+x_n\mathbf{c}_n=\mathbf{b}$.

$\implies$ $\mathbf{b}\in \operatorname{span} \{\mathbf{c}_1,\mathbf{c}_2,\ldots,\mathbf{c}_n\}$.

$\implies$ $\operatorname{span}\{\underbrace{\mathbf{c}_1,\ldots,\mathbf{c}_n}_{\text{columns of $A$}}\}=\operatorname{span}\{\underbrace{\mathbf{c}_1,\ldots,\mathbf{c}_n,\mathbf{b}}_{\textcolor{red}{\text{columns of ($A\mid\mathbf{b}$)}}}\}$.

$\implies$ $\operatorname{im}(\textcolor{red}{A})=\operatorname{im}(\textcolor{red}{A\mid\mathbf{b}})$,

$\implies$ $\operatorname{rank}(A)=\operatorname{rank}(A\mid\mathbf{b})$.

\medskip

$(\Leftarrow)$
Conversely, suppose $\operatorname{rank}(A)=\operatorname{rank}(A\mid\mathbf{b})$.

$\implies$ $\operatorname{im}(A)\subseteq\operatorname{im}(A\mid\mathbf{b})$.

Since these two subspaces have the same dimension, they must be equal

$\implies$ $\operatorname{im}(A)=\operatorname{im}(A\mid\mathbf{b})$.

Thus, $\mathbf{b}\in\operatorname{span}\{\mathbf{c}_1,\ldots,\mathbf{c}_n\}$.

$\implies$ $\exists$ $x_1,\ldots,x_n$ s.t. $x_1\mathbf{c}_1+\cdots+x_n\mathbf{c}_n=\mathbf{b}$.
Equivalently, there $\exists$ $\mathbf{x}$ s.t. $A\mathbf{x}=\mathbf{b}$.

Therefore, the system is consistent.

\end{proof}

\textbf{Theorem 7.4.4. (Priciple of linearity)}
Suppose $\mathbf{v}$ is a solution to the system of equations
$(A\mid\mathbf{b})$. Then the general solution of $(A\mid\mathbf{b})$
is given by
\[
\mathbf{v}+\ker(A)
\coloneqq
\left\{
\mathbf{v}+\mathbf{n}:\mathbf{n}\in\ker(A)
\right\}.
\]

\begin{proof}
Given that $\mathbf{v}$ is a solution of $(A\mid\mathbf{b})$, we have $A\mathbf{v}=\mathbf{b}$ $\textup{  (1)}$

Suppose $\mathbf{u}$ is another solution of $(A\mid\mathbf{b})$. Then $A\mathbf{u}=\mathbf{b}$ $\textup{  (2)}$

Consider \textup{(2)} $-$ \textup{(1)}:
\begin{align*}
A\mathbf{u}-A\mathbf{v}
    &=\mathbf{b}-\mathbf{b}\\
\Longleftrightarrow\quad
A(\mathbf{u}-\mathbf{v})
    &=\mathbf{0}\\
\Longleftrightarrow\quad
\mathbf{u}-\mathbf{v}
    &\in\ker(A)\\
\Longleftrightarrow\quad
\mathbf{u}-\mathbf{v}
    &=\mathbf{n}
    \quad\text{for some }\mathbf{n}\in\ker(A)\\
\Longleftrightarrow\quad
\mathbf{u}
    &=\mathbf{v}+\mathbf{n}
    \quad\text{for some }\mathbf{n}\in\ker(A).
\end{align*}

Hence, the general solution of $(A\mid\mathbf{b})$ is $\mathbf{v}+\ker(A)$.

\end{proof}
\begin{Remark}
In particular, the system $(A\mid\mathbf{b})$ has a \textbf{unique} solution iff.
\[
\begin{cases}
\operatorname{rank}(A)=\operatorname{rank}(A\mid\mathbf{b}),\\[0.3em]
\ker(A)=\{\mathbf{0}\}.
\end{cases}
\]
\end{Remark}

\subsection{Example}

\subsubsection{General systems: Spanning}
Does $(3,3,2)$ belong to $\operatorname{span}\{(1,1,1),(1,-1,0)\}$?
\begin{solution}
Suppose it does. Then there $\exists$ $a,b\in\mathbb{R}$ s.t.
\[
(3,3,2)=a(1,1,1)+b(1,-1,0).
\]

Equivalently,
\[
\begin{cases}
a+b=3,\\
a-b=3,\\
a=2.
\end{cases}
\]

Consider the coefficient matrix
\[
A=
\begin{pmatrix}
1 & 1\\
1 & -1\\
1 & 0
\end{pmatrix}.
\]
We have $\operatorname{rank}(A)=2$.

On the other hand, the augmented matrix is
\[
(A\mid\mathbf{b})=
\left(
\begin{array}{cc|c}
1 & 1  & 3\\
1 & -1 & 3\\
1 & 0  & 2
\end{array}
\right),
\]
and $\operatorname{rank}(A\mid\mathbf{b})=3$.

Thus, $\operatorname{rank}(A)\neq\operatorname{rank}(A\mid\mathbf{b})$,
so the corresponding system is inconsistent.

Hence,
$\boxed{
(3,3,2)
\notin
\operatorname{span}\{(1,1,1),(1,-1,0)\}
}.$

\end{solution}

\subsubsection{General systems: Linear Transformations}

Consider the linear transformation $T:\mathbb{R}[x]_{\leq 2}\to\mathbb{R}[x]_{\leq 2}$ defined by
\[
T(a+bx+cx^2)=(a+b+c)+(a-b+c)x+(a+c)x^2.
\]
Does $3+3x+2x^2$ belong to $\operatorname{im}(T)$?

\begin{solution}
Suppose it does. $\implies$ $3+3x+2x^2=T(a+bx+cx^2)$ for some $a,b,c\in\mathbb{R}$,

By comparing the coefficients, we obtain
\[
\begin{cases}
a+b+c=3,\\
a-b+c=3,\\
a+c=2.
\end{cases}
\tag{*}
\]

The coefficient matrix of $(*)$ is
\[
A=
\begin{pmatrix}
1 & 1  & 1\\
1 & -1 & 1\\
1 & 0  & 1
\end{pmatrix},
\]
and $\operatorname{rank}(A)=2$.

On the other hand,
\[
\operatorname{rank}
\left(
\begin{array}{ccc|c}
1 & 1  & 1 & 3\\
1 & -1 & 1 & 3\\
1 & 0  & 1 & 2
\end{array}
\right)
=3.
\]

Therefore, $\operatorname{rank}(A)\neq\operatorname{rank}(A\mid\mathbf{b})$, so $(*)$ has no solution.

Hence,
$\boxed{
3+3x+2x^2\notin\operatorname{im}(T)
}.$

\end{solution}

\subsubsection{Inhomogeneous systems}
Consider the inhomogeneous system
$\begin{cases}
x_1+2x_2+3x_3=4,\\
5x_1+6x_2+7x_3=8.
\end{cases}
$ Find the general solution.

\begin{solution}
The general solution is given by
\[
\left\{
\begin{pmatrix}
x_1\\
x_2\\
x_3
\end{pmatrix}
=
\begin{pmatrix}
-2+t\\
3-2t\\
t
\end{pmatrix}
:t\in\mathbb{R}
\right\}.
\]
\begin{align*}
&=
\left\{
\begin{pmatrix}
-2\\
3\\
0
\end{pmatrix}
+
t
\begin{pmatrix}
1\\
-2\\
1
\end{pmatrix}
:t\in\mathbb{R}
\right\}\\
&=
\begin{pmatrix}
-2\\
3\\
0
\end{pmatrix}
+
\operatorname{span}
\left\{
\begin{pmatrix}
1\\
-2\\
1
\end{pmatrix}
\right\}.
\end{align*}

Here,
$\begin{pmatrix}
-2\\
3\\
0
\end{pmatrix}$ is a particular solution, while
$\operatorname{span}
\left\{
\begin{pmatrix}
1\\
-2\\
1
\end{pmatrix}
\right\}
=\ker
\begin{pmatrix}
1&2&3\\
5&6&7
\end{pmatrix}$.


Thus, the general solution has the form 
$\boxed{
\text{particular solution}+\ker(A)
}$.
\end{solution}

\section{TFAE Theorem Ver.1}

\subsection{Theorem}

\textbf{Theorem 7.5.1.}
Let $T:V\to W$ be a linear transformation s.t.
\[
\dim(V)=\dim(W)<\infty.
\]
Therefore, having fixed bases, $T$ can be represented by a \textbf{square matrix} $A$.

The following are equivalent (TFAE).
\begin{enumerate}
    \item $T$ is injective.
    
    \item $\ker(T)=\{\mathbf{0}_V\}$.
    
    \item $\operatorname{nullity}(T)=0$.
    
    \item $\operatorname{rank}(T)=\dim(W)$
    ($T$ has full rank).
    
    \item $\operatorname{im}(T)=W$.
    
    \item $T$ is surjective.
    
    \item $T$ is an isomorphism (bijective linear map).
    
    \item $A$ is invertible ($A^{-1}$ exists).
    
    \item $\ker(A)=\{\mathbf{0}\}$.
    
    \item $A$ has full rank.
    
    \item All the rows of $A$ are linearly independent.
    
    \item All the columns of $A$ are linearly independent.
    
    \item The system of equations $A\mathbf{x}=\mathbf{b}$
    has a unique solution, namely $\mathbf{x}=A^{-1}\mathbf{b}$.
\end{enumerate}

\newpage

\subsection{Example}

\subsubsection{Prove that a matrix is invertible}
Let $A\in M_n(F)$. It is known that $A^k=\mathbf{0}$
for some positive integer $k$.

Prove that $A+I_n$ is invertible.

\begin{proof}
Consider $\ker(A+I_n)$.

Let $\mathbf{v}\in\ker(A+I_n)$. Then
\begin{align*}
(A+I_n)\mathbf{v}&=\mathbf{0}\\
A\mathbf{v}+\mathbf{v}&=\mathbf{0}\\
A\mathbf{v}&=-\mathbf{v}.
\end{align*}

Multiplying $A$ to the left on both sides gives
\[
A^2\mathbf{v}
=-A\mathbf{v}
=(-1)(-\mathbf{v})
=(-1)^2\mathbf{v}.
\]

Inductively, $\textcolor{red}{\mathbf{0} =}$ $A^k\mathbf{v}=(-1)^k\mathbf{v}\implies \mathbf{v}=\mathbf{0}$.
Hence, $\ker(A+I_n)=\{\mathbf{0}\}$.

By TFAE, since $A+I_n$ is a square matrix, $A+I_n$ is invertible.
\end{proof}

\section{Change of coordinate matrix}

\subsection{Definition}

\textbf{Definition 7.6.1.}
Let $V$ be a finite-dimensional vector space. Suppose $S_1$ and $S_2$
are two bases of $V$. Then the \emph{change of coordinate matrix},
or \emph{transition matrix}, from $S_1$-basis to $S_2$-basis is the matrix
$[\operatorname{id}_V]_{S_1}^{S_2}$.

\begin{Remark}
Let $Q=[\operatorname{id}_V]_{S_1}^{S_2}$.

\begin{enumerate}
    \item Since $\operatorname{id}_V$ is an isomorphism, the matrix $Q$
    is invertible. Moreover,
    \[
    Q^{-1}
    =
    [\operatorname{id}_V]_{S_2}^{S_1}.
    \]

    \item For any $\mathbf{v}\in V$, we have
    \[
    [\operatorname{id}_V]_{S_1}^{S_2}
    [\mathbf{v}]_{S_1}
    =
    [\mathbf{v}]_{S_2}.
    \]

    In other words, the matrix
    $[\operatorname{id}_V]_{S_1}^{S_2}$ converts the coordinate vector
    of $\mathbf{v}$ in the $S_1$-basis into the coordinate vector of
    $\mathbf{v}$ in the $S_2$-basis. This explains the name
    ``change of coordinate'' matrix.
\end{enumerate}
\end{Remark}

\textbf{Definition 7.6.2.}
Let $V$ be a vector space. A linear transformation $T:V\to V$
is called a \emph{linear operator on $V$} or an \emph{endomorphism on $V$}.

\subsection{Theorem}

\textbf{Theorem 7.6.1.}
Let $V$ be finite-dimensional and let $T:V\to V$ be a linear operator.
Suppose
\begin{itemize}
    \item $S_1,S_2$ are two bases of the vector space $V$;
    
    \item $Q=[\operatorname{id}_V]_{S_1}^{S_2}$
    is the change of coordinate matrix from $S_1$ to $S_2$.
\end{itemize}

Then we have an equality of matrices:
\[
\boxed{[T]_{S_1}^{S_1}=Q^{-1}[T]_{S_2}^{S_2}Q}.
\]

\begin{proof}
Consider the following diagram:
\[
\begin{tikzcd}[row sep=large, column sep=huge]
V_{S_1}
    \arrow[r,
        "{\textcolor{red}{[\textcolor{black}{T}]_{S_1}^{S_1}}}"]
    \arrow[d,
        "{\textcolor{red}{
        Q=[\textcolor{black}{\operatorname{id}_V}]_{S_1}^{S_2}
        }}"']
&
V_{S_1}
\\
V_{S_2}
    \arrow[r,
        "{\textcolor{red}{[\textcolor{black}{T}]_{S_2}^{S_2}}}"']
&
V_{S_2}
\arrow[u,
        "{\textcolor{red}{
        [\textcolor{black}{\operatorname{id}_V}]_{S_2}^{S_1}=Q^{-1}
        }}"']
\end{tikzcd}
\]

Here the subscripts indicate the bases used for the coordinates.
Since
\[
\textcolor{red}{
Q=[\textcolor{black}{\operatorname{id}_V}]_{S_1}^{S_2}
}
\qquad\text{and}\qquad
\textcolor{red}{
[\textcolor{black}{\operatorname{id}_V}]_{S_2}^{S_1}=Q^{-1},
}
\]
So $\textcolor{red}{[\textcolor{black}{T}]_{S_1}^{S_1}=Q^{-1}[\textcolor{black}{T}]_{S_2}^{S_2}Q.}$

\end{proof}

\subsection{Example}

\subsubsection{Change of coordinate matrices}
Let $B$ be the standard basis of $\mathbb{R}^2$ and
be the basis $S=\{(1,3),(2,5)\}$ of $\mathbb{R}^2$.

\begin{enumerate}
    \item[(a)] Find the change of coordinate matrix from $S$ to $B$.
    
    \item[(b)] Find the change of coordinate matrix from $B$ to $S$.
    
    \item[(c)] Find $[(-2,3)]_S$.
\end{enumerate}

\newpage

\begin{solution}
Let $B=\{\mathbf{e}_1,\mathbf{e}_2\}$.

\begin{enumerate}
    \item[(a)] We have
    \begin{align*}
        \operatorname{id}(\textcolor{red}{1,3})
        &= (1,3)
         = \mathbf{e}_1+3\mathbf{e}_2,\\
        \operatorname{id}(\textcolor{red}{2,5})
        &= (2,5)
         = 2\mathbf{e}_1+5\mathbf{e}_2.
    \end{align*}
    Therefore,
    $[\operatorname{id}]_{\textcolor{red}{S}}^B=
    \begin{pmatrix}
        1 & 2\\
        3 & 5
    \end{pmatrix}.$

    \item[(b)] By the inverse relation between change of coordinate matrices,
    \begin{align*}
        [\operatorname{id}]_B^{\textcolor{red}{S}}
        =
        \left([\operatorname{id}]_{\textcolor{red}{S}}^B\right)^{-1}
        =
        \begin{pmatrix}
            1 & 2\\
            3 & 5
        \end{pmatrix}^{-1}
        =
        \begin{pmatrix}
            -5 & 2\\
            3 & -1
        \end{pmatrix}.
    \end{align*}

    \item[(c)] Since $[(-2,3)]_B
    =
    \begin{pmatrix}
        -2\\
        3
    \end{pmatrix}$,
    we have
    \begin{align*}
        [(-2,3)]_{\textcolor{red}{S}}
        =
        [\operatorname{id}]_B^{\textcolor{red}{S}}[(-2,3)]_B
        =
        \begin{pmatrix}
            -5 & 2\\
            3 & -1
        \end{pmatrix}
        \begin{pmatrix}
            -2\\
            3
        \end{pmatrix}
        =
        \begin{pmatrix}
            16\\
            -9
        \end{pmatrix}.
    \end{align*}

    Indeed,
    \[
    16\textcolor{red}{(1,3)}-9\textcolor{red}{(2,5)}=\textcolor{red}{(-2,3)}.
    \]
\end{enumerate}
\end{solution}

\subsubsection{Matrix Representations under Different Bases}
Let $T:\mathbb{R}^2\to\mathbb{R}^2$ be the linear operator on $\mathbb{R}^2$ defined by reflection by the 

line $y=-x$. Let
\[
B=\{(1,0),(0,1)\}
\qquad\text{and}\qquad
S=\{(1,1),(1,-1)\}
\]
be two bases of $\mathbb{R}^2$.

\begin{enumerate}
    \item[(a)] Find the matrix $X=[T]_B^B$.

    \item[(b)] Find the matrix $Y=[T]_S^S$.

    \item[(c)] Find the change of coordinate matrix $Q$ from $S$ to $B$.

    \item[(d)] Write down the relation between the matrices $X,Y,Q$.
\end{enumerate}

\begin{solution}
\begin{enumerate}
\item[(a)]
\begin{align*}
T(1,0)
    &= (0,-1)
     = 0\mathbf{e}_1-\mathbf{e}_2,\\
T(0,1)
    &= (-1,0)
     = -\mathbf{e}_1+0\mathbf{e}_2.
\end{align*}

Therefore, $X=[T]_B^B=
\begin{pmatrix}
0 & -1\\
-1 & 0
\end{pmatrix}$.

\item[(b)]
\begin{align*}
T(\textcolor{orange}{(1,1)})
    &=(-1,-1)=-1\textcolor{orange}{(1,1)}+0\textcolor{orange}{(1,-1)},\\
T(\textcolor{orange}{(1,-1)})
    &=(1,-1)=0\textcolor{orange}{(1,1)}+1\textcolor{orange}{(1,-1)}.
\end{align*}

Therefore, $Y=[T]_{\textcolor{orange}{S}}^{\textcolor{orange}{S}}=
\begin{pmatrix}
-1 & 0\\
0 & 1
\end{pmatrix}$.

\item[(c)]
\begin{align*}
\operatorname{id}(\textcolor{orange}{(1,1)})
    &=(1,1)
     =\mathbf{e}_1+\mathbf{e}_2,\\
\operatorname{id}(\textcolor{orange}{(1,-1)})
    &=(1,-1)
     =\mathbf{e}_1-\mathbf{e}_2.
\end{align*}

Hence, $Q=[\operatorname{id}]_{\textcolor{orange}{S}}^B=
\begin{pmatrix}
1 & 1\\
1 & -1
\end{pmatrix}$.

\item[(d)]
Consider the following diagram:
\[
\begin{tikzcd}[row sep=large, column sep=huge]
\mathbb{R}^2_{\textcolor{orange}{S\text{-basis}}}
    \arrow[r,
        "{\textcolor{orange}{
        \textcolor{red}{Y=[T]}_S^S
        }}"]
    \arrow[d,
        "{\textcolor{red}{
        Q=[\operatorname{id}]_{\textcolor{orange}{S}}^{\textcolor{blue}{B}}
        }}"']
&
\mathbb{R}^2_{\textcolor{orange}{S\text{-basis}}}
\\
\mathbb{R}^2_{\textcolor{blue}{B\text{-basis}}}
    \arrow[r,
        "{\textcolor{blue}{
        \textcolor{red}{X=[T]}_B^B
        }}"']
&
\mathbb{R}^2_{\textcolor{blue}{B\text{-basis}}}
    \arrow[u,
        "{\textcolor{red}{
        [\operatorname{id}]_{\textcolor{blue}{B}}^{\textcolor{orange}{S}}=Q^{-1}
        }}"']
\end{tikzcd}
\]

Since
\[
\textcolor{blue}{X=[T]_B^B},
\qquad
\textcolor{red}{Y=[T]_S^S},
\]
following the diagram gives
\[
\boxed{
\textcolor{red}{Y}
=
Q^{-1}\textcolor{blue}{X}Q
}.
\]

Equivalently,
\[
\boxed{
\textcolor{blue}{X}
=
Q\textcolor{red}{Y}Q^{-1}
}.
\]

\end{enumerate}
\end{solution}

\end{document}